%%%%%ENTREES%%%%
\begin{tabarg}
\multicolumn{3}{l}{\titretab {Entr\'ees}}\\
\hline
$L_x$ & Coeff., \texttt{float}& Coordonn�e du point de mesure de la vitesse 
de l'air par rapport au mobile, sur l'axe X du rep�re dans lequel
sont exprim�es les mesures de la vitesse par rapport au sol (m).\\
$L_y$ & Coeff., \texttt{float}& Coordonn�e du point de mesure de la vitesse 
de l'air par rapport au mobile, sur l'axe Y du rep�re dans lequel
sont exprim�es les mesures de la vitesse par rapport au sol (m).\\
$L_z$ & Coeff., \texttt{float}& Coordonn�e du point de mesure de la vitesse 
de l'air par rapport au mobile, sur l'axe Z du rep�re dans lequel
sont exprim�es les mesures de la vitesse par rapport au sol (m).\\

$V_p$ & Vecteur, \texttt{float} & Vitesse propre ($m/s$) \\
$\alpha$ & Vecteur, \texttt{float} & Angle d'attaque, positif lorsque le vent 
relatif monte (vient de plus bas que la plate-forme) ($radians$) \\
$\beta$ & Vecteur, \texttt{float} & Angle de d\'erapage, positif lorsque le 
vent relatif vient de tribord ($radians$)\\

$V_x$ & Vecteur, \texttt{float} & Premi�re composante de la vitesse du 
mobile par rapport au sol (en g�n�ral, suivant l'axe longitudinal de
la plate-forme, positif vers l'avant) ($m/s$)\\
$V_y$ & Vecteur, \texttt{float} & Seconde composante de la vitesse du 
mobile par rapport au sol (en g�n�ral, suivant l'axe 
transversal de la plate-forme, positif de b�bord vers tribord) ($m/s$)\\ 
$V_z$ & Vecteur, \texttt{float} & Troisi�me composante de la vitesse du
mobile par rapport au sol (en g�n�ral, suivant l'axe 
"vertical" de la plate-forme, positif vers le bas) ($m/s$)\\

$Vr_x$ & Vecteur, \texttt{float} & Vitesse de rotation suivant l'axe X, positif 
quand $\phi$ cro\^{\i}t($rad/s$) \\
$Vr_y$ & Vecteur, \texttt{float} & Vitesse de rotation suivant l'axe Y, positif
quand $\theta$ cro\^{\i}t($rad/s$) \\
$Vr_z$ & Vecteur, \texttt{float} & Vitesse de rotation suivant l'axe Z, positif
quand $\psi$ cro\^{\i}t ($rad/s$) \\

$\theta$ & Vecteur, \texttt{float} & Angle de tangage du mobile, positif quand
le mobile est cabr\'e ($radians$)\\
$\phi$ & Vecteur, \texttt{float} & Angle de roulis du mobile, positif quand
le cot� tribord est enfonc\'e ($radians$)\\
$\psi$ & Vecteur, \texttt{float} & Cap g\'eographique du mobile ($radians$) \\
\end{tabarg}

\medskip
%%%%%SORTIES%%%%
\begin{tabarg}
\multicolumn{3}{l}{\titretab {Sorties}}\\
\hline
$U$ & Vecteur, \texttt{float} & Composante zonale du vent (Ouest en Est) ($m/s$)\\
$V$ & Vecteur, \texttt{float} & Composante m\'eridienne du vent (Sud au Nord) ($m/s$)\\
$W$ & Vecteur, \texttt{float} & Composante verticale du vent (Bas en Haut) ($m/s$)\\
\end{tabarg}
